\section{Tweaks and Balancing}

While the core mechanics of the game might be perfectly functional, the game can only reach its full potential if there
isn't just one winning strategy every time. This is where balancing comes into play, balancing is the act of tweaking
the game so that certain gameplay styles are not too powerful or too weak.

\subsection{Economic Balancing}

In the Economic Simulation chapter I mentioned that the preferences of people is based on samples of a normal
distribution along with some weights and shifts: this is where they come into play. One of the first things I noticed
when playing the game is that high prices were overpowered: I you are just able to sell one product for 17 quadrillion
Franks you you beat everybody else trying to play tactfully. One of the early things I had to add is a hard price cap.
Another key aspect is to lower the size of the market. In earlier versions of the game, the market had around 100'000
population. Because companies start very small and the market is very big, this lead to virtually no competition, as
there was always a huge shortage of products.

\subsection{Research and Development}

Research is one of the hardest mechanics to optimise while keeping the game fun. The issue is that research having too
high an effectiveness would defeat the need to ever meaningfully change anything in your company, on the other side, if
research is too weak then what's the point? The solution I found to this problem was to slightly lower the effectiveness
of research and only apply it new products, this way players are forced to develop new products and maybe try out some
different things compared to what they had before.
